% 第九卷 8.4.2 节：四选一题集上的可靠性图（教学示例，非实验数据）。
% 100道题按置信度分入四个等宽区间，各区间样本数为20、30、30、20；
% 区间平均置信度为0.35、0.55、0.75、0.95，实际正确率为0.40、0.50、0.60、0.80；
% 按式 eq:nlp-eval-ece 计算 ECE=0.2*0.05+0.3*0.05+0.3*0.15+0.2*0.15=0.10。
\begin{tikzpicture}[figure picture]
  \begin{axis}[
    width=\linewidth,
    height=58mm,
    xmin=0.2,xmax=1,
    ymin=0,ymax=1,
    xlabel={置信度},
    ylabel={实际正确率},
    xtick={0.2,0.4,0.6,0.8,1.0},
    ytick={0,0.2,0.4,0.6,0.8,1.0},
    grid=major,
    legend style={at={(0.5,1.03)},anchor=south,legend columns=3,
      /tikz/every even column/.append style={column sep=4mm}},
  ]
    \addplot[ybar interval,fill=FigureBlueFill,draw=FigureBlue,line width=.6pt]
      coordinates {(0.2,0.40) (0.4,0.50) (0.6,0.60) (0.8,0.80) (1.0,0.80)};
    \addlegendentry{区间正确率}
    \addplot[FigureStroke,densely dashed,line width=.8pt,domain=0.2:1,samples=2] {x};
    \addlegendentry{完全校准}
    \addplot[only marks,mark=*,mark size=1.6pt,FigureRed]
      coordinates {(0.35,0.40) (0.55,0.50) (0.75,0.60) (0.95,0.80)};
    \addlegendentry{（平均置信度，正确率）}
    \draw[figure link,FigureRed,line width=.9pt] (axis cs:0.35,0.35) -- (axis cs:0.35,0.40);
    \draw[figure link,FigureRed,line width=.9pt] (axis cs:0.55,0.50) -- (axis cs:0.55,0.55);
    \draw[figure link,FigureRed,line width=.9pt] (axis cs:0.75,0.60) -- (axis cs:0.75,0.75);
    \draw[figure link,FigureRed,line width=.9pt] (axis cs:0.95,0.80) -- (axis cs:0.95,0.95);
    \node[figure note,anchor=south] at (axis cs:0.3,0.03) {$n_b=20$};
    \node[figure note,anchor=south] at (axis cs:0.5,0.03) {$n_b=30$};
    \node[figure note,anchor=south] at (axis cs:0.7,0.03) {$n_b=30$};
    \node[figure note,anchor=south] at (axis cs:0.9,0.03) {$n_b=20$};
  \end{axis}
\end{tikzpicture}
